% compendium/execution-trading-costs/impact-models.tex
% Entries migrated on 2026-10-15 from the legacy liquidity-and-market-impact blocks (slice 7);
% dispositions and source checks in notes/migration.md, the blocks verbatim in
% notes/artefacts/market-making-lob/derivations.tex.

\section{Price impact}

\begin{qfdefinition}[Price impact: unperturbed, observed and trade prices; reaction, prediction, transient and permanent impact]\label{def:price-impact}
\qffield{Statement} \emph{Price impact} is the effect of trading on prices: trades cause price moves that would not otherwise have happened \citep[chap.~1, sec.~1.1.2, p.~5]{wide-webster-price-impact}. Webster separates three prices for a trader whose trades $Q$ are part of all trades $L$ \citep[chap.~2, Definition~2.2.6, pp.~53--54]{wide-webster-price-impact}:
\begin{itemize}
\item the \emph{unperturbed price} $S_t$, the price without the trades $L$ and their impact;
\item the \emph{perturbed} or observed price $P_t=S_t+I_t$, where $I_t$ is the impact of $L$;
\item the \emph{trade price} $\widetilde P_{t-}=S_{t-}+I_{t-}+s_{t-}$, where $s$ is the instantaneous transaction (spread) cost of the trader's own fill, and the subscript $t-$ marks that a fill at $t$ executes against the state just before $t$.
\end{itemize}
Bouchaud, Bonart, Donier and Gould separate three impacts of a market order executed at $t$, at horizon $\ell$, with $m$ the mid-price and $\mathcal F_t$ everything that may have prompted the order \citep[chap.~11, sec.~11.2, eqs.~(11.1)--(11.4), pp.~210--211]{catalogue-0528}:
\[ I^{\mathrm{react}}=\mathbb E[m_{t+\ell}\mid\mathrm{exec}_t,\mathcal F_t]-\mathbb E[m_{t+\ell}\mid\text{no exec}_t,\mathcal F_t],\qquad I^{\mathrm{obs}}=\mathbb E[m_{t+\ell}\mid\mathrm{exec}_t]-m_t, \]
\[ I^{\mathrm{pred}}=\mathbb E[m_{t+\ell}\mid\text{no exec}_t,\mathcal F_t]-m_t,\qquad I^{\mathrm{obs}}=I^{\mathrm{react}}+I^{\mathrm{pred}}. \]
Over time, impact has a \emph{transient} part that decays and a \emph{permanent} part. The permanent part combines reaction and prediction: it is higher for traders with predictive power than for noise trades \citep[chap.~12, sec.~12.4, p.~240]{catalogue-0528}.

\qffield{Notation} $I^{\mathrm{react}}$, $I^{\mathrm{obs}}$, $I^{\mathrm{pred}}$ are the reaction, observed and prediction impacts. The source writes $I^{\mathrm{react.}}_{t+\ell}(\mathrm{exec}_t\mid\mathcal F_t)$ and so on; the time and conditioning arguments are dropped here. The identity $I^{\mathrm{obs}}=I^{\mathrm{react}}+I^{\mathrm{pred}}$ needs every expectation conditioned on $\mathcal F_t$; as printed, (11.2) conditions on execution alone, and the authors state that the conditioning on $\mathcal F_t$ is always implicit \citep[same sec., p.~210]{catalogue-0528} (Interpretation: without it the two sides differ by $\mathbb E[m_{t+\ell}\mid\mathrm{exec}_t]-\mathbb E[m_{t+\ell}\mid\mathrm{exec}_t,\mathcal F_t]$). Webster's $P$ is a price, not the $P$ of other entries.

\qffield{Conditions} Two readings of impact frame the literature \citep[chap.~11, sec.~11.1, pp.~208--210]{catalogue-0528}. In the efficient-market reading, traders forecast prices and trade accordingly, so trades and later price changes correlate even if trades move nothing; in Hasbrouck's words, quoted there, orders forecast prices. In the order-flow reading, trading itself moves prices, and impact exists even in a model with no information at all, such as the Santa Fe model. Reality probably lies between. If the mid-price were a martingale, observed and reaction impact would coincide; they do not, because the information behind the order raises the expected price even without the trade (p.~211). Studies with random or cash-driven trades find that, over short horizons, mechanical impact is about the same for random and allegedly informed trades; prediction shows only at longer horizons (p.~212).

\qffield{Practice} Reaction impact can be estimated by trading at random, by randomising execution of signals, or by conditioning on the rest of the market's order-sign imbalance; prediction impact is very hard to estimate \citep[chap.~11, sec.~11.2, pp.~211--212]{catalogue-0528}. The counterfactual cannot be observed: a trader never sees both outcomes of one trade, which Webster calls the causality challenge \citep[chap.~1, sec.~1.1.3, pp.~6--7]{wide-webster-price-impact}. Webster traces impact modelling through four communities: economists' stylised models (Kyle), mathematical finance's control problems (Almgren--Chriss, Obizhaeva--Wang, Gatheral), econophysics' empirical response functions and propagators, and books from quantitative trading desks \citep[chap.~1, sec.~1.1.5, pp.~9--15]{wide-webster-price-impact}.

\qffield{Pitfalls} ``Temporary impact'' is ambiguous. Almgren and Chriss call the instantaneous transaction cost ``temporary market impact'' \citep[chap.~2, n.~3, p.~53]{wide-webster-price-impact}, a cost paid on one's own fill, whereas transient impact is a price displacement that decays. The legacy definition (``the price concession needed to obtain immediacy \dots expected to decay as liquidity replenishes'') mixes the two; say which is meant (Interpretation). Calling a long-run move permanent impact needs care. Long-horizon impact is hard to measure because price variance grows with the horizon. Correlated metaorders, the same decision continued over days, can look like permanent impact. Separating mechanical from informational impact is difficult \citep[chap.~12, sec.~12.4, p.~240]{catalogue-0528}. The legacy caveat also listed market and sector moves to remove. The passages checked suggest conditioning on the rest of the market's order-sign imbalance (p.~212) but do not discuss market or sector returns. The legacy taxonomy of three lenses (theoretical, reduced-form, practitioner) is a simplification of Webster's four communities.

\qffield{See also} \hyperref[mod:propagator]{propagator model}; \hyperref[res:no-price-manipulation]{absence of price manipulation}; \hyperref[def:price-formation-process]{price formation process} (temporary and permanent impact by consensus, Lehalle and Laruelle); \hyperref[res:spread-component-identification]{spread component identification} (transitory and permanent components of the spread); \hyperref[mod:kyle-single-auction]{Kyle single-auction model} (permanent impact from information); \hyperref[mod:grossman-miller]{Grossman--Miller model} (a temporary concession for immediacy).

\qffield{Further reading} The legacy blocks it replaces, verbatim: \nolinkurl{notes/artefacts/market-making-lob/derivations.tex}.

\qffield{Index} price impact; market impact; unperturbed price; perturbed price; trade price; instantaneous transaction cost; reaction impact; observed impact; prediction impact; transient impact; permanent impact; temporary impact; causality challenge; metaorders; price discovery and price formation

\qffield{Status} Verification: partly checked 2026-10-15, against the text layers of local copies of two books (Webster's handbook, \texttt{wide-webster-price-impact}, CRC 2023 first edition, pp.~3--15 and 53--54, and Bouchaud et al., \texttt{catalogue-0528}, Cambridge 2018 first printing, ISBN 9781107156050, pp.~208--212 and 240; the page image of p.~211 for eqs.~(11.2)--(11.4)). Stability: draft. Currency: timeless (concepts); the empirical statements as of 2018.
\end{qfdefinition}

\begin{qfmodel}[Propagator model: transient impact of past trades]\label{mod:propagator}
\qffield{Statement} In market-order time, with $\varepsilon_n=\pm1$ the sign of the $n$th market order and $\xi_n$ price changes not induced by trades, the mid-price is \citep[chap.~13, sec.~13.2, eq.~(13.7), p.~252, page image]{catalogue-0528}
\[ m_t=m_{t_0}+\sum_{t_0\le n<t}G(t-n)\,\varepsilon_n+\sum_{t_0\le n<t}\xi_n. \]
The \emph{propagator} $G(\ell)$ is the impact at lag $\ell$ of one trade, expected to decrease with $\ell$; $G_\infty=\lim_{\ell\to\infty}G(\ell)$ is the permanent reaction impact. Returns then split into immediate impact, the decay of past impacts and noise, $r_t=G_1\varepsilon_t+\sum_{n<t}K(t-n)\varepsilon_n+\xi_t$ with $K(\ell)=G(\ell+1)-G(\ell)$ (same sec., eq.~(13.8)). Equation~(13.7) defines the family of linear transient impact models (TIM).

\qffield{Notation} $m_t$ mid-price just before event $t$; $r_t=m_{t+1}-m_t$; $G_1=G(1)$; $C(\ell)=\mathbb E[\varepsilon_t\varepsilon_{t+\ell}]$ the sign autocorrelation; $R(\ell)=\mathbb E[\varepsilon_t(m_{t+\ell}-m_t)]$ the response function (observed impact).

\qffield{Conditions} Why impact must be transient: if each trade had constant permanent impact $G_1$, then $\mathbb E[r_tr_{t+\ell}]=G_1^2C(\ell)$, and the long memory of order signs would make returns strongly autocorrelated. That contradicts the nearly flat signature plots of real markets, so much of each trade's impact must relax \citep[chap.~13, secs.~13.1--13.2, pp.~249--251]{catalogue-0528}. The simple model assumes impact independent of trade volume, time and market conditions, which the authors call gross oversimplifications, relaxed in their chap.~14 \citep[same chap., p.~252]{catalogue-0528}.

\qffield{Practice} The observed response and the bare propagator differ when signs are correlated: $R(\ell)=G(\ell)+\sum_{0<n<\ell}G(\ell-n)C(n)+\sum_{n>0}[G(\ell+n)-G(n)]C(n)$, which reduces to $G(\ell)$ only for uncorrelated signs. Estimating $G$ from $R$ and $C$ through this relation is sensitive to finite-size effects \citep[same chap., sec.~13.2.1, eq.~(13.10), p.~253]{catalogue-0528}. Webster calls the Obizhaeva--Wang model a simple form of the propagator model \citep[chap.~1, p.~10]{wide-webster-price-impact}; in its generalized form impact follows $dI_t=-\beta_tI_t\,dt+\lambda_t\,dQ_t$ \citep[chap.~2, Definition~2.8.1, p.~97]{wide-webster-price-impact}. Interpretation: that is an exponentially decaying kernel in calendar time, $G\propto e^{-\beta\ell}$ for constant $\beta$. The model originates with Bouchaud, Gefen, Potters and Wyart (2004), quoted by Webster \citep[chap.~1, p.~5]{wide-webster-price-impact} and not consulted here (\texttt{wide-bouchaud-propagator}).

\qffield{Pitfalls} The legacy block described the price as the sum of signed impulses ``with different decay profiles''. In the simple model every trade shares one propagator $G$; trade-dependent profiles need the generalised models of chap.~14, which were not read. Reading the measured response $R(\ell)$ as the impact of a single trade ignores the sign autocorrelation. Interpretation: Hasbrouck's generalized Roll model, in \hyperref[res:spread-component-identification]{spread component identification}, is the constant-propagator case applied to the efficient price with independent signs: the trade moves it permanently by $\lambda$.

\qffield{See also} \hyperref[def:price-impact]{price impact}; \hyperref[res:no-price-manipulation]{absence of price manipulation} (restrictions on $G$ and on the Obizhaeva--Wang parameters); \hyperref[def:intraday-overnight-returns]{returns by time scale}; \hyperref[def:volatility-components]{fundamental and transitory volatility} (the variance ratio and transient components).

\qffield{Further reading} The original paper, \citep{wide-bouchaud-propagator}, not consulted. The legacy block it replaces, verbatim: \nolinkurl{notes/artefacts/market-making-lob/derivations.tex}.

\qffield{Index} propagator; transient impact model; TIM; impact decay; response function; order-sign autocorrelation; long memory of order flow; signature plot; market-order time; Obizhaeva--Wang model; resilience

\qffield{Status} Verification: partly checked 2026-10-15 (against the text layer of a local copy of \texttt{catalogue-0528}, Cambridge 2018 first printing, pp.~249--253, with the page image of p.~252; \texttt{wide-webster-price-impact}, CRC 2023, pp.~5, 10 and 97). Stability: draft. Currency: timeless (model); the empirical motivation as of 2018.
\end{qfmodel}

\begin{qfresult}[Absence of price manipulation in impact models]\label{res:no-price-manipulation}
\qffield{Statement} In an impact model, \emph{price manipulation} is a round trip, a strategy that starts and ends with no position, whose expected cost is negative. Webster quotes this definition from Gatheral (2010), which was not consulted \citep[chap.~1, n.~3, p.~7]{wide-webster-price-impact}. A usable model must exclude it. Bouchaud, Bonart, Donier and Gould write the expected cost of a closed trajectory with trading rate $j$ and price $x^*$ as $\mathcal C=\int_0^Tj(u)\,x^*(u)\,du$ with $\int_0^Tj(u)\,du=0$. They show that in their latent-liquidity model it is a non-negative quadratic form, so $\mathcal C\ge0$ for every schedule, generalising a result for the linear propagator model \citep[chap.~19, sec.~19.5, pp.~361--362, page image]{catalogue-0528}. In Webster's generalized Obizhaeva--Wang model, buying $\Delta$ in one jump at $t$ and selling it at $t'>t$, without a spread, earns
\[ (S_{t'}-S_t)\Delta-\tfrac12\lambda_t(2\beta_t+\gamma'_t)(t'-t)\Delta^2+o(t'-t), \]
so if $2\beta+\gamma'<0$, large, quick round trips earn arbitrary expected profits from the trader's own impact. The no-manipulation condition is $2\beta+\gamma'>0$, after Fruth, Schöneborn and Urusov (2013), not consulted \citep[chap.~2, sec.~2.7.1, Example~2.7.1, p.~94, page image]{wide-webster-price-impact}.

\qffield{Notation} $S$ unperturbed price; $\Delta$ shares; impact $I$ with $dI=-\beta I\,dt+\lambda\,dQ$, decay $\beta$ per unit time, push $\lambda=e^\gamma$ in price per share, $\gamma'$ its rate of change. As Webster's profit expression implies, the trader's buy executes at $S_t+\tfrac12\lambda_t\Delta$ on average, the sell at $S_{t'}+e^{-\int_t^{t'}\beta_s\,ds}\lambda_t\Delta-\tfrac12\lambda_{t'}\Delta$.

\qffield{Conditions} Webster's example assumes no spread, no prior impact, jump trades and exact unwinding; the expected profit from $S$ is zero only if $S$ is a martingale. The condition $2\beta+\gamma'>0$ rules out only the paradoxes that yield arbitrary profits. It does not rule out transaction-triggered manipulation, such as selling first to depress the price for later buying \citep[same sec., pp.~94--95]{wide-webster-price-impact}. In a square-root model it bounds forecast volume growth, $v'_t/v_t<4\beta_t$: trading activity cannot double in less than a quarter of the impact half-life, so rising volume curves can create manipulation opportunities while falling ones cannot \citep[same chap., sec.~2.7.3, Example~2.7.2, p.~96]{wide-webster-price-impact}.

\qffield{Practice} Using an impact model that admits profitable closed trajectories in execution algorithms leads to instabilities or counter-intuitive trades \citep[chap.~19, p.~362]{catalogue-0528}, and Webster asks that a model rule out manipulation paradoxes \citep[chap.~1, sec.~1.1.4, p.~8]{wide-webster-price-impact}. Interpretation, the accounting behind any check: a trader who buys and sells $Q$ shares at average execution prices $\bar P_B$ and $\bar P_S$, impact included, and pays explicit costs $F$ in total, ends with $Q(\bar P_S-\bar P_B)-F$ and no position, hence no mark-to-market term. With a per-share cost $c$ on each side, $F=2Qc$; with $c$ quoted per round-trip share, $F=Qc$. Spread counted in $\bar P$ must not be counted again in $c$.

\qffield{Pitfalls} A realised profitable round trip is not price manipulation: the profit may come from $S_{t'}-S_t$, from information, or from luck. Manipulation concerns expected profit under a model, with no information (Interpretation). Impact that treats buys and sells alike per unit is not by itself a no-manipulation guarantee. Webster's example is linear and symmetric trade by trade, yet admits manipulation when liquidity changes faster than impact decays (Interpretation). The legacy principle ``impact depends on net pressure and risk, so a zero-net sequence cannot mechanically create a higher average sale price'' holds in Harris's instantaneous, permanent-impact example but is not a general result. The legacy proposition stated as a test only that identity, with $c$ undefined per side or per round trip. Strategic bluffing, which fools other traders through their inferences, is a different mechanism from manipulation through one's own impact in a model, where nobody is fooled (\hyperref[def:bluffing-manipulation]{bluffing and manipulation}).

\qffield{See also} \hyperref[mod:propagator]{propagator model}; \hyperref[def:price-impact]{price impact}; \hyperref[def:bluffing-manipulation]{bluffing and manipulation} (Harris's asymmetric-impact bluff).

\qffield{Further reading} Gatheral, ``No-Dynamic-Arbitrage and Market Impact'' \citep{catalogue-0501}, a required reading of \texttt{price-impact-execution}, not consulted. The legacy block it replaces, verbatim: \nolinkurl{notes/artefacts/market-making-lob/derivations.tex}.

\qffield{Index} price manipulation; no-dynamic-arbitrage; round trip; closed trajectory; Obizhaeva--Wang model; transaction-triggered manipulation; volume curves; impact decay; liquidity; execution cost

\qffield{Status} Verification: partly checked 2026-10-15 (against \texttt{wide-webster-price-impact}, CRC 2023, pp.~7--8 and 93--97, with the page image of p.~94; and \texttt{catalogue-0528}, Cambridge 2018 first printing, pp.~361--362, with its page image; the example rederived independently). Stability: draft. Currency: timeless (theory).
\end{qfresult}
